\chapter{Stan wiedzy}
\label{ch:sota}

Rozdział ten przedstawia tło literaturowe pracy w~czterech obszarach: estymacji pozy z~wideo monocular 2D, biomechaniki biegu, klasyfikacji faz biegu i~chodu metodami uczenia maszynowego oraz walidacji metod 2D względem trójwymiarowego motion capture. Na tej podstawie identyfikuję lukę badawczą, którą praca wypełnia. Szczegóły metod zastosowanych w~niniejszej pracy opisuję w~rozdziale~\ref{ch:materialy-metody}; tutaj koncentruję się na dorobku innych autorów.

\section{Estymacja pozy z~wideo monocular 2D}
\label{sec:sota-pose}

Klasyczna analiza ruchu opiera się na systemach motion capture z~markerami (np.~Vicon, Qualisys), które oferują wysoką dokładność, lecz są drogie, wymagają laboratorium i~doświadczonego operatora. Rozwój metod estymacji pozy bez markerów (\textit{markerless pose estimation}) otworzył możliwość analizy ruchu z~nagrań ze~zwykłej kamery. Frameworki takie jak OpenPose~\cite{cao2021openpose} oraz MediaPipe BlazePose~\cite{bazarevsky2020blazepose}, trenowane m.in.\ na zbiorze COCO~\cite{lin2014coco}, wykrywają keypointy ciała w~każdej klatce wideo, działając na pojedynczej kamerze (\textit{monocular}) i~obrazie dwuwymiarowym.

Roggio i~wsp.~\cite{roggio2024review} w~przeglądzie dziewięciu modeli estymacji pozy wskazują na ich potencjał w~zastosowaniach klinicznych, sportowych i~ergonomicznych jako tania, nieinwazyjna alternatywa dla systemów markerowych, zaznaczając jednocześnie wyzwania związane z~dokładnością, jakością danych i~zasłonięciami. Edriss i~wsp.~\cite{frontiers2025review} w~przeglądzie poświęconym analizie ruchu w~sporcie podkreślają, że rozwiązania 2D są ekonomiczne i~proste w~użyciu, lecz wciąż mają trudność z~rejestracją ruchów poza płaszczyzną obrazu (\textit{out-of-plane}). Porównanie konkretnych standardów keypointów oraz uzasadnienie wyboru MediaPipe Pose w~niniejszej pracy przedstawiam w~sekcji~\ref{sec:mediapipe}.

\section{Biomechanika biegu --- współczynniki i~ich znaczenie}
\label{sec:sota-wspolczynniki}

Cykl biegowy różni się od cyklu chodu: w~biegu występuje faza lotu (obie stopy w~powietrzu), a~zanika faza podwójnego podparcia charakterystyczna dla chodu~\cite{novacheck1998}. Z~tego powodu progi referencyjne i~wzorce ruchu ustalone dla chodu klinicznego nie przenoszą się wprost na bieg sportowy.

Literatura biomechaniczna dostarcza ugruntowanych współczynników opisujących technikę biegu. Heiderscheit i~wsp.~\cite{heiderscheit2011} wykazali, że manipulacja kadencją wpływa na obciążenia stawów i~wiąże się z~overstridingiem oraz kątem kolana w~momencie kontaktu. Souza~\cite{souza2016} omawia czas kontaktu z~podłożem i~wzorzec lądowania stopy w~kontekście analizy wideo. Novacheck~\cite{novacheck1998} opisuje fazy cyklu biegowego, pochylenie tułowia, kąty stawów i~oscylację pionową. Daoud i~wsp.~\cite{daoud2012} powiązali wzorzec lądowania stopy z~częstością kontuzji u~biegaczy długodystansowych, a~Robinson i~wsp.~\cite{robinson1987} zaproponowali wskaźnik symetrii (\textit{Symmetry Index}) do ilościowego opisu różnic między stronami ciała. Przeglądy nowsze, jak Figueiredo i~wsp.~\cite{figueiredo2025cadence}, potwierdzają związek kadencji z~prewencją kontuzji. Konkretne zakresy referencyjne tych współczynników i~sposób ich obliczania w~niniejszej pracy zestawiam w~sekcji~\ref{sec:obliczanie-wspolczynnikow}.

\section{Klasyfikacja faz biegu i~chodu metodami uczenia maszynowego}
\label{sec:sota-klasyfikatory}

Rozpoznawanie faz cyklu ruchu jest przedmiotem wielu prac, najczęściej opartych na czujnikach inercyjnych (IMU) noszonych na ciele. Vu i~wsp.~\cite{vu2020review} dokonali przeglądu algorytmów detekcji faz chodu dla protez kończyn dolnych. Su i~wsp.~\cite{su2020dcnn} zastosowali głęboką sieć konwolucyjną z~danych IMU, a~Jeon i~Lee~\cite{jeon2024bilstm} dwukierunkową sieć LSTM (BiLSTM) odporną na zmiany kierunku ruchu --- architekturę pokrewną zastosowanej w~niniejszej pracy.

W~kontekście biegu klasyczne metody uczenia maszynowego wykorzystywano m.in.\ do wykrywania zmian zmęczeniowych: Dimmick i~wsp.~\cite{dimmick2023fatigue} zastosowali las losowy (Random Forest) do analizy biomechaniki biegu, a~Martínez-Gramage i~wsp.~\cite{martinez2020rf} --- do redukcji ryzyka kontuzji u~młodych triathlonistów. Prace łączące estymację pozy z~wideo z~klasyfikacją uczenia maszynowego pojawiają się rzadziej: Ali i~wsp.~\cite{ali2024dmd} połączyli estymację pozy (m.in.\ MediaPipe) z~klasyfikatorem SVM oraz LSTM-CNN do rozróżniania patologii chodu, a~Young i~wsp.~\cite{young2023smartphone} zaproponowali ocenę biegu z~pojedynczej kamery smartfona. Większość prac dotyczy jednak chodu klinicznego i~opiera się na czujnikach noszonych, nie na czystym wideo.

\section{Walidacja estymacji 2D względem 3D motion capture}
\label{sec:sota-walidacja}

Kluczowym wątkiem literatury jest pytanie, na ile estymacja pozy z~wideo 2D dorównuje referencyjnym systemom 3D. Stenum i~wsp.~\cite{stenum2021} porównali OpenPose z~systemem 3D MoCap dla chodu, uzyskując błąd średni rzędu 0{,}02~s dla metryk czasowych oraz kilku stopni dla kątów stawów. Wskazali przy tym istotne ograniczenie: długość kroku jest estymowana najdokładniej, gdy osoba znajduje się w~centrum pola widzenia kamery --- a~więc wynik zależy od pozycji w~kadrze. Menychtas i~wsp.~\cite{menychtas2023} porównali OpenPose i~MediaPipe z~systemem Vicon, raportując błędy w~ruchach prostopadłych do płaszczyzny obrazu oraz problemy z~lokalizacją keypointu kostki. Praca oparta na metodzie HGcnMLP~\cite{hgcnmlp2024} pokazała wysoką zgodność estymacji 3D ze~smartfona z~systemem Vicon, lecz dane zbierano wyłącznie przy jednej orientacji kamery ($90^\circ$), bez walidacji innych perspektyw.

Z~prac tych wyłania się spójny obraz ograniczeń estymacji 2D: błędy dla ruchów poza płaszczyzną obrazu, zależność od perspektywy oraz od pozycji w~polu widzenia. Ograniczenia te są w~literaturze \textit{identyfikowane}, ale rzadko \textit{systematycznie analizowane} pod kątem wpływu na konkretne współczynniki --- co stanowi część luki opisanej poniżej.

\section{Luka badawcza}
\label{sec:sota-luka}

Z~przeglądu literatury wynika, że żadna z~analizowanych prac nie łączy jednocześnie następujących elementów, które wyznaczają zakres niniejszej pracy:

\begin{enumerate}
    \item \textbf{Ukierunkowanie na bieg, nie chód.} Analizowane prace dotyczą niemal wyłącznie chodu (gait), zwykle w~kontekście klinicznym. Bieg sportowy na bieżni ma odmienne fazy (występuje faza lotu), inne progi kadencji i~czasu kontaktu oraz inne wartości referencyjne --- stanowi więc odrębną niszę.
    \item \textbf{Systematyczne porównanie czterech klasyfikatorów na małym zbiorze.} Większość prac proponuje jeden model. W~niniejszej pracy porównuję cztery konfiguracje (las losowy na cechach surowych i~inżynierowanych oraz dwie sieci BiLSTM), traktując mały, dobrze opisany zbiór jako świadomą specyfikację, nie wadę.
    \item \textbf{Korekta proporcji obrazu jako krok preprocessingu.} Menychtas i~wsp.~\cite{menychtas2023} identyfikują błędy dla ruchów poza płaszczyzną obrazu, lecz nie raportują badania wpływu korekty proporcji kadru na jakość klasyfikatora faz.
    \item \textbf{Silnik reguł rekomendacji oparty na literaturze.} Prace cytują wartości referencyjne (Heiderscheit, Souza, Novacheck, Robinson, Daoud), ale nie operacjonalizują ich w~postaci wykonywalnych reguł z~poziomami ważności, źródłami i~regułami łączonymi.
    \item \textbf{Automatyczna detekcja niskiej wiarygodności pomiaru.} Stenum~\cite{stenum2021} i~Menychtas~\cite{menychtas2023} identyfikują problem perspektywy, lecz żadna z~prac nie sygnalizuje automatycznie przypadków, w~których dany współczynnik staje się nieinterpretowalny.
\end{enumerate}

Pełne sformułowanie pytań badawczych i~hipotez, wynikających z~tej luki, przedstawiam w~rozdziale~\ref{ch:wstep}.
